% !TEX program = xelatex
% Unified CRPS-uncertainty summary across the three experiments.
% Table generated by make_crps_table.R -> crps_table_body.tex.
\documentclass[12pt]{article}

\usepackage{fontspec}
\usepackage{xeCJK}
\setCJKmainfont[AutoFakeSlant=.2, AutoFakeBold=3]{AR PL KaitiM Big5}

\usepackage[margin=1in]{geometry}
\usepackage{amsmath, amssymb}
\usepackage{booktabs}
\usepackage{array}
\usepackage{hyperref}
\hypersetup{colorlinks=true, linkcolor=blue!45!black, urlcolor=blue!45!black}

\title{\textbf{量化 CRPS 的不確定性：三個實驗的統一結果}}
\author{plain CRPS $+$ 站層 bootstrap／複製層 Wilcoxon}
\date{2026 年 7 月}

\begin{document}
\maketitle

\section*{可以，而且和 Brier 是同一件事}

「量化 CRPS 的不確定性」與「量化 Brier 的不確定性」是\textbf{同一個問題}：把每個
預測化約成一個 per-unit 分數後，後面的重抽／配對檢定機器完全不變。換的只是分數
算法。事實上兩者有精確關係
\[
  \mathrm{CRPS}(F,y) \;=\; \int_{-\infty}^{\infty}
  \underbrace{\bigl(F(L) - \mathbf{1}\{y\le L\}\bigr)^2}_{=\,\mathrm{Brier}_L}\, dL ,
\]
你的門檻-$L$ Brier 就是 CRPS 被積函數在 $L$ 的值。CRPS 把它對所有門檻積分，所以
每個預測攜帶的資訊更多、變異更小、\textbf{檢定力更高}。

三個 arm 的 CRPS 不確定性都算好了，用各自對口的獨立單位：臭氧用\textbf{站層
bootstrap}，模擬與 time-block 用 $10$ 組資料的\textbf{複製層配對檢定}。

\bigskip
% Auto-generated by make_crps_table.R -- do not edit by hand.
\begin{table}[htbp]
\centering
\caption{CRPS-difference uncertainty across the three arms (whole-
distribution proper score; $\Delta = \mathrm{CRPS}_{\text{first}} -
\mathrm{CRPS}_{\text{second}}$, negative favours the temporal/extension
model). Ozone uses a $95\%$ site-clustered bootstrap CI ($p$ two-sided);
simulation and time-block use a $95\%$ paired-$t$ CI over the $10$ data sets
($p_W$ one-sided Wilcoxon, H$_1$: temporal better). The reading follows each
row's own test at $\alpha=0.05$; the two-sided time-block CIs graze $0$ even
where the one-sided $p_W$ is significant. CRPS reveals what $q_{95}$ Brier could not:
AR(2) is (mildly) \emph{worse} for real-data interpolation yet \emph{better}
for time-block extrapolation.}
\label{tab:crps-uncertainty}
\small
\begin{tabular}{@{}llccc@{}}
\toprule
Arm & Contrast & $\Delta$ CRPS [95\% CI] & $p$ & reading \\
\midrule
\multicolumn{5}{@{}l}{\emph{Interpolation --- spatial hold-out}}\\
Ozone 200/200 & AR(2) $-$ AR(1) & $+0.021$ {\scriptsize $[+0.003,+0.04]$} & $p{=}0.026$ & \textbf{worse} \\
Ozone 200/200 & MRTS $-$ none & $-0.017$ {\scriptsize $[-0.053,+0.019]$} & $p{=}0.344$ & ns \\
Ozone 300/100 & AR(2) $-$ AR(1) & $+0.039$ {\scriptsize $[+0.0037,+0.072]$} & $p{=}0.027$ & \textbf{worse} \\
Sim K{=}1 strong & AR(2) $-$ AR(1) & $-0.000083$ {\scriptsize $[-0.00052,+0.00036]$} & $p_W{=}0.342$ & ns \\
Sim K{=}1 weak & AR(2) $-$ AR(1) & $-0.00000097$ {\scriptsize $[-0.00023,+0.00022]$} & $p_W{=}0.581$ & ns \\
Sim K{=}5 strong & AR(2) $-$ AR(1) & $+0.00052$ {\scriptsize $[-0.0069,+0.0079]$} & $p_W{=}0.581$ & ns \\
Sim K{=}5 weak & AR(2) $-$ AR(1) & $-0.0039$ {\scriptsize $[-0.013,+0.0052]$} & $p_W{=}0.270$ & ns \\
\addlinespace
\multicolumn{5}{@{}l}{\emph{Extrapolation --- time-block forecast}}\\
Near-unit-root (1--5) & AR(2) $-$ iid & $-0.8$ {\scriptsize $[-1.7,+0.11]$} & $p_W{=}0.026$ & \textbf{better} \\
Near-unit-root (1--5) & AR(2) $-$ AR(1) & $-0.66$ {\scriptsize $[-1.5,+0.21]$} & $p_W{=}0.033$ & \textbf{better} \\
\bottomrule
\end{tabular}
\end{table}


\section*{CRPS 揭露了 Brier 看不到的事}

\paragraph{內插（臭氧、模擬空間 hold-out）：AR(2) 不只無益，在真實資料上甚至微幅有害。}
在 $q_{95}$ Brier 上，臭氧的 AR(2)$-$AR(1) 差異 CI 涵蓋 $0$（測不出）。改用 CRPS
後，AR(2) 相對 AR(1) 的 CRPS \textbf{顯著較差}，且在 $200/200$ 與 $300/100$ 兩個
split 都成立（$p\approx0.03$）。效果雖小（CRPS $\approx 5$ ppb，差 $0.02$--$0.04$），
但方向一致、可複現：多加第二階輕微拖累\emph{整條}預測分布。MRTS 則在 CRPS 上
仍與基準無法區分。（\emph{兩點修正，見 folder 06/07/08}：此「較差」既屬
\textbf{bulk} 效應——尾部的 Brier 與 twCRPS 皆 ns；又對站間空間相依\textbf{不穩健}
——空間 block bootstrap 下 $p$ 由 $0.03$ 升到 $0.08$--$0.11$、CI 含 $0$。故無論
對口與否或統計穩健與否，它都不承載結論。）

\paragraph{外推（time-block forecast）：plain CRPS 較佳，但（修正後）是 bulk。}
同一個模型換到向前預測的任務，AR(2) 相對 iid 與 AR(1) 的 \emph{plain} CRPS 改善
在近單位根設定下達（單尾）顯著（$p_W\approx0.03$）。\textbf{但 folder 07 的尾部帶
twCRPS 顯示此增益同樣在 bulk}：一約束到尾部，$\Delta$ 掉回 $\approx0$
（$p_W\approx0.46$），與單門檻 Brier 一致。AR(2) 改善的是未來日的\emph{水準}，
不是超標機率。

\paragraph{合起來（修正後）：尾部一律測不到。}folder 06/07 用尾部帶 twCRPS 把
兩邊的 plain-CRPS 訊號都打回 bulk。用論文真正在乎的\textbf{尾部}分數（單門檻
Brier 與檢定力更高的 twCRPS）衡量：
\[
  \text{內插：尾部 ns} \qquad\text{且}\qquad \text{外推：尾部 ns.}
\]
也就是說，\textbf{AR(2) 在所有情境都不改變門檻超標的預測技巧}；plain CRPS 的差異
（臭氧較差、time-block 較佳）都是分布主體（level）效應，與超標目標無關。這比
原本「AR(2) 是外推工具」的敘述更保守，卻由\emph{兩個}尾部分數、跨\emph{三個}
arm 一致支撐，證據力更強。

\section*{兩點但書}

\begin{enumerate}
  \item \textbf{plain CRPS vs threshold-weighted CRPS。}臭氧模型是 censored（只配
        尾部）。plain CRPS 對整條分布等權、故被 bulk 主導；對齊「超標預測」目標的是
        twCRPS $=\int \mathrm{Brier}_L\,w(L)\,dL$（尾部加權）。\textbf{folder 06/07 已
        建好 twCRPS} 並以它做出上述 bulk-vs-tail 修正；量化方式與 plain CRPS 完全
        相同，只換權重 $w(L)$。
  \item \textbf{多重性。}跨 arm 已累積數個邊界星號；正式寫入論文時應做多重比較
        校正（審查 \#2），對最極端／最稀疏的顯著格保守解讀。
\end{enumerate}

\medskip
\noindent\footnotesize 程式：\verb|ozone/US-all-auto/crps_bootstrap.R|、
\verb|simstudy/ar2_crps_export.R|、
\verb|block1_positive_control/timeblock_crps_export.R|。

\end{document}
